% =============================================================================
% appendix.tex — the essential appendix, \input by paper.tex after \appendix.
%   Appendix A: the model (via % =============================================================================
% model.tex — A simple model of delivery-mode choice and physician time.
% % =============================================================================
% model.tex — A simple model of delivery-mode choice and physician time.
% % =============================================================================
% model.tex — A simple model of delivery-mode choice and physician time.
% \input{model} from paper.tex. Delivers the five predictions tested in the
% empirical sections (P1-P5) and the policy corollary.
% =============================================================================

\section{A simple model of delivery-mode choice and physician time}
\label{sec:model}

\subsection{Setup}

A physician attends a pregnancy whose labor, absent intervention, begins at a
random time $\tilde{t}$ distributed (approximately) uniformly over the days and
hours of the week. Clinically, the onset of spontaneous labor does not respect
the calendar, which is the identifying assumption behind our timing tests. The
physician chooses between attending a vaginal delivery ($v$) and performing a
\emph{prelabor scheduled} cesarean ($c$) at a time $t$ of her choosing; if labor
starts before a cesarean is performed, the delivery may still convert to an
\emph{in-labor} cesarean for clinical reasons.

\paragraph{Fees.} The physician receives $f_c$ for a cesarean and $f_v$ for a
vaginal delivery, where $f_v$ includes the separately-billed hourly labor
assistance, the \emph{economic} vaginal fee. In our data $f_c - f_v \le 0$ in
the largest states.

\paragraph{Time.} A cesarean takes a fixed block $\tau_c$ ($\approx$1 hour) at
the chosen time $t$. A vaginal delivery requires attending labor for a random
duration $\tau_v$ (billed mean $\approx$2.7 hours, with $E[\tau_v] > \tau_c$) at
the random time $\tilde{t}$. The opportunity cost of physician time is
\[
w(t) \;=\;
\begin{cases}
w_0 & t \in \text{weekday business hours}, \\
w_0 + \Delta & t \in \text{nights, weekends, holidays},
\end{cases}
\qquad \Delta > 0 .
\]
Committing to an event of unknown timing additionally carries an on-call cost
$\kappa \ge 0$: the physician must hold open calendar slots she could otherwise
fill (an option cost of unschedulable time).

\paragraph{Agency.} A vaginal delivery yields the patient a net clinical benefit
$\theta \sim F$ (with $\theta < 0$ when a cesarean is medically indicated). The
physician is an imperfect agent and internalizes $\alpha\theta$, $\alpha \in
(0,1]$.

\subsection{The delivery-mode decision}

The physician schedules a prelabor cesarean if and only if
\begin{equation}
\underbrace{f_c - f_v}_{\text{fee gap } \le 0}
\;+\;
\underbrace{\Pi}_{\text{convenience premium}}
\;>\; \alpha\,\theta,
\qquad
\Pi \;\equiv\; E\!\left[w(\tilde{t})\right] E[\tau_v] + \kappa - w_0\,\tau_c \;>\; 0,
\label{eq:choice}
\end{equation}
so the cutoff is $\theta^{*} = \left(f_c - f_v + \Pi\right)/\alpha$ and the
cesarean share is $F(\theta^{*})$. The convenience premium $\Pi$ is strictly
positive even when the fee gap is negative: labor falls disproportionately on
high-$w$ times ($E[w(\tilde{t})] > w_0$), lasts longer ($E[\tau_v] > \tau_c$),
and blocks the calendar ($\kappa$). Conditional on choosing $c$, the physician
schedules it at $\arg\min_t w(t)$ (weekday business hours) and, to
guarantee the procedure precedes $\tilde{t}$ (whose hazard rises steeply after
gestational week 39), books it at weeks 37--39, earlier than spontaneous labor
would deliver.

\subsection{Predictions and the corresponding tests}

\begin{itemize}
\item[\textbf{P1.}] \emph{Fees are not the operative margin.} The cesarean share
responds to the fee gap only through $f(\theta^{*})/\alpha$; when $\Pi$
dominates the wedge in \eqref{eq:choice}, the discretionary mass is already
inframarginal and fee (semi-)elasticities are $\approx 0$.
\hfill [\emph{Tested:} fee-gap regressions and fee-swing falsification, Tables
\ref{tab:not_a_price} and \ref{tab:fee_shock}; even a court order to triple
vaginal pay changed nothing.]
\item[\textbf{P2.}] \emph{The scheduling fingerprint.} Prelabor cesareans bunch
at $\arg\min_t w(t)$ (weekday business hours) and fall on weekends and
holidays; in-labor cesareans inherit $\tilde{t}$'s uniform timing and weakly
\emph{rise} on weekends (marginal-$\theta$ women who labor before being
scheduled convert). \hfill [\emph{Tested:} day-of-week, holiday and
hour-of-birth clustering, and the prelabor/in-labor split: Figures
\ref{fig:dow}--\ref{fig:hours}, Tables \ref{tab:scheduling} and
\ref{tab:mechanism_checks}.]
\item[\textbf{P3.}] \emph{Scarcity.} $\Delta$ and $\kappa$ rise where
obstetricians per birth are scarce (fewer substitutes to cover on-call time), so
weekend dips are larger there. \hfill [\emph{Tested:} density interaction, Table
\ref{tab:heterogeneity}.]
\item[\textbf{P4.}] \emph{The early-term cost.} Guaranteeing $t < \tilde{t}$
requires booking before the week-39 hazard spike, shifting prelabor cesareans to
37--38 weeks and producing excess early-term birth. \hfill [\emph{Tested:}
gestational-age distributions and early-term gap, Figure \ref{fig:gestation},
Table \ref{tab:health}.]
\item[\textbf{P5.}] \emph{Supply, not demand.} If the timing preference were the
mother's, the weekend dip would concentrate among mothers with the strongest
demand for scheduling (proxied by schooling); it does not. \hfill
[\emph{Tested:} education interaction, Table \ref{tab:heterogeneity}.]
\end{itemize}

\subsection{Policy corollary}

Interventions on $f$, such as equalizing or even tripling the vaginal fee, leave
$\theta^{*}$ nearly unchanged because the operative wedge is $\Pi$, not the fee
gap. The model thereby rationalizes our fee nulls jointly with the timing
evidence, and directs attention to organizational levers that reduce $\Pi$
itself: on-call/laborist shift coverage (which lowers $\kappa$ and detaches
$w(\tilde t)$ from the attending physician's calendar), midwife-led continuity
of care, and hospital scheduling norms.
 from paper.tex. Delivers the five predictions tested in the
% empirical sections (P1-P5) and the policy corollary.
% =============================================================================

\section{A simple model of delivery-mode choice and physician time}
\label{sec:model}

\subsection{Setup}

A physician attends a pregnancy whose labor, absent intervention, begins at a
random time $\tilde{t}$ distributed (approximately) uniformly over the days and
hours of the week. Clinically, the onset of spontaneous labor does not respect
the calendar, which is the identifying assumption behind our timing tests. The
physician chooses between attending a vaginal delivery ($v$) and performing a
\emph{prelabor scheduled} cesarean ($c$) at a time $t$ of her choosing; if labor
starts before a cesarean is performed, the delivery may still convert to an
\emph{in-labor} cesarean for clinical reasons.

\paragraph{Fees.} The physician receives $f_c$ for a cesarean and $f_v$ for a
vaginal delivery, where $f_v$ includes the separately-billed hourly labor
assistance, the \emph{economic} vaginal fee. In our data $f_c - f_v \le 0$ in
the largest states.

\paragraph{Time.} A cesarean takes a fixed block $\tau_c$ ($\approx$1 hour) at
the chosen time $t$. A vaginal delivery requires attending labor for a random
duration $\tau_v$ (billed mean $\approx$2.7 hours, with $E[\tau_v] > \tau_c$) at
the random time $\tilde{t}$. The opportunity cost of physician time is
\[
w(t) \;=\;
\begin{cases}
w_0 & t \in \text{weekday business hours}, \\
w_0 + \Delta & t \in \text{nights, weekends, holidays},
\end{cases}
\qquad \Delta > 0 .
\]
Committing to an event of unknown timing additionally carries an on-call cost
$\kappa \ge 0$: the physician must hold open calendar slots she could otherwise
fill (an option cost of unschedulable time).

\paragraph{Agency.} A vaginal delivery yields the patient a net clinical benefit
$\theta \sim F$ (with $\theta < 0$ when a cesarean is medically indicated). The
physician is an imperfect agent and internalizes $\alpha\theta$, $\alpha \in
(0,1]$.

\subsection{The delivery-mode decision}

The physician schedules a prelabor cesarean if and only if
\begin{equation}
\underbrace{f_c - f_v}_{\text{fee gap } \le 0}
\;+\;
\underbrace{\Pi}_{\text{convenience premium}}
\;>\; \alpha\,\theta,
\qquad
\Pi \;\equiv\; E\!\left[w(\tilde{t})\right] E[\tau_v] + \kappa - w_0\,\tau_c \;>\; 0,
\label{eq:choice}
\end{equation}
so the cutoff is $\theta^{*} = \left(f_c - f_v + \Pi\right)/\alpha$ and the
cesarean share is $F(\theta^{*})$. The convenience premium $\Pi$ is strictly
positive even when the fee gap is negative: labor falls disproportionately on
high-$w$ times ($E[w(\tilde{t})] > w_0$), lasts longer ($E[\tau_v] > \tau_c$),
and blocks the calendar ($\kappa$). Conditional on choosing $c$, the physician
schedules it at $\arg\min_t w(t)$ (weekday business hours) and, to
guarantee the procedure precedes $\tilde{t}$ (whose hazard rises steeply after
gestational week 39), books it at weeks 37--39, earlier than spontaneous labor
would deliver.

\subsection{Predictions and the corresponding tests}

\begin{itemize}
\item[\textbf{P1.}] \emph{Fees are not the operative margin.} The cesarean share
responds to the fee gap only through $f(\theta^{*})/\alpha$; when $\Pi$
dominates the wedge in \eqref{eq:choice}, the discretionary mass is already
inframarginal and fee (semi-)elasticities are $\approx 0$.
\hfill [\emph{Tested:} fee-gap regressions and fee-swing falsification, Tables
\ref{tab:not_a_price} and \ref{tab:fee_shock}; even a court order to triple
vaginal pay changed nothing.]
\item[\textbf{P2.}] \emph{The scheduling fingerprint.} Prelabor cesareans bunch
at $\arg\min_t w(t)$ (weekday business hours) and fall on weekends and
holidays; in-labor cesareans inherit $\tilde{t}$'s uniform timing and weakly
\emph{rise} on weekends (marginal-$\theta$ women who labor before being
scheduled convert). \hfill [\emph{Tested:} day-of-week, holiday and
hour-of-birth clustering, and the prelabor/in-labor split: Figures
\ref{fig:dow}--\ref{fig:hours}, Tables \ref{tab:scheduling} and
\ref{tab:mechanism_checks}.]
\item[\textbf{P3.}] \emph{Scarcity.} $\Delta$ and $\kappa$ rise where
obstetricians per birth are scarce (fewer substitutes to cover on-call time), so
weekend dips are larger there. \hfill [\emph{Tested:} density interaction, Table
\ref{tab:heterogeneity}.]
\item[\textbf{P4.}] \emph{The early-term cost.} Guaranteeing $t < \tilde{t}$
requires booking before the week-39 hazard spike, shifting prelabor cesareans to
37--38 weeks and producing excess early-term birth. \hfill [\emph{Tested:}
gestational-age distributions and early-term gap, Figure \ref{fig:gestation},
Table \ref{tab:health}.]
\item[\textbf{P5.}] \emph{Supply, not demand.} If the timing preference were the
mother's, the weekend dip would concentrate among mothers with the strongest
demand for scheduling (proxied by schooling); it does not. \hfill
[\emph{Tested:} education interaction, Table \ref{tab:heterogeneity}.]
\end{itemize}

\subsection{Policy corollary}

Interventions on $f$, such as equalizing or even tripling the vaginal fee, leave
$\theta^{*}$ nearly unchanged because the operative wedge is $\Pi$, not the fee
gap. The model thereby rationalizes our fee nulls jointly with the timing
evidence, and directs attention to organizational levers that reduce $\Pi$
itself: on-call/laborist shift coverage (which lowers $\kappa$ and detaches
$w(\tilde t)$ from the attending physician's calendar), midwife-led continuity
of care, and hospital scheduling norms.
 from paper.tex. Delivers the five predictions tested in the
% empirical sections (P1-P5) and the policy corollary.
% =============================================================================

\section{A simple model of delivery-mode choice and physician time}
\label{sec:model}

\subsection{Setup}

A physician attends a pregnancy whose labor, absent intervention, begins at a
random time $\tilde{t}$ distributed (approximately) uniformly over the days and
hours of the week. Clinically, the onset of spontaneous labor does not respect
the calendar, which is the identifying assumption behind our timing tests. The
physician chooses between attending a vaginal delivery ($v$) and performing a
\emph{prelabor scheduled} cesarean ($c$) at a time $t$ of her choosing; if labor
starts before a cesarean is performed, the delivery may still convert to an
\emph{in-labor} cesarean for clinical reasons.

\paragraph{Fees.} The physician receives $f_c$ for a cesarean and $f_v$ for a
vaginal delivery, where $f_v$ includes the separately-billed hourly labor
assistance, the \emph{economic} vaginal fee. In our data $f_c - f_v \le 0$ in
the largest states.

\paragraph{Time.} A cesarean takes a fixed block $\tau_c$ ($\approx$1 hour) at
the chosen time $t$. A vaginal delivery requires attending labor for a random
duration $\tau_v$ (billed mean $\approx$2.7 hours, with $E[\tau_v] > \tau_c$) at
the random time $\tilde{t}$. The opportunity cost of physician time is
\[
w(t) \;=\;
\begin{cases}
w_0 & t \in \text{weekday business hours}, \\
w_0 + \Delta & t \in \text{nights, weekends, holidays},
\end{cases}
\qquad \Delta > 0 .
\]
Committing to an event of unknown timing additionally carries an on-call cost
$\kappa \ge 0$: the physician must hold open calendar slots she could otherwise
fill (an option cost of unschedulable time).

\paragraph{Agency.} A vaginal delivery yields the patient a net clinical benefit
$\theta \sim F$ (with $\theta < 0$ when a cesarean is medically indicated). The
physician is an imperfect agent and internalizes $\alpha\theta$, $\alpha \in
(0,1]$.

\subsection{The delivery-mode decision}

The physician schedules a prelabor cesarean if and only if
\begin{equation}
\underbrace{f_c - f_v}_{\text{fee gap } \le 0}
\;+\;
\underbrace{\Pi}_{\text{convenience premium}}
\;>\; \alpha\,\theta,
\qquad
\Pi \;\equiv\; E\!\left[w(\tilde{t})\right] E[\tau_v] + \kappa - w_0\,\tau_c \;>\; 0,
\label{eq:choice}
\end{equation}
so the cutoff is $\theta^{*} = \left(f_c - f_v + \Pi\right)/\alpha$ and the
cesarean share is $F(\theta^{*})$. The convenience premium $\Pi$ is strictly
positive even when the fee gap is negative: labor falls disproportionately on
high-$w$ times ($E[w(\tilde{t})] > w_0$), lasts longer ($E[\tau_v] > \tau_c$),
and blocks the calendar ($\kappa$). Conditional on choosing $c$, the physician
schedules it at $\arg\min_t w(t)$ (weekday business hours) and, to
guarantee the procedure precedes $\tilde{t}$ (whose hazard rises steeply after
gestational week 39), books it at weeks 37--39, earlier than spontaneous labor
would deliver.

\subsection{Predictions and the corresponding tests}

\begin{itemize}
\item[\textbf{P1.}] \emph{Fees are not the operative margin.} The cesarean share
responds to the fee gap only through $f(\theta^{*})/\alpha$; when $\Pi$
dominates the wedge in \eqref{eq:choice}, the discretionary mass is already
inframarginal and fee (semi-)elasticities are $\approx 0$.
\hfill [\emph{Tested:} fee-gap regressions and fee-swing falsification, Tables
\ref{tab:not_a_price} and \ref{tab:fee_shock}; even a court order to triple
vaginal pay changed nothing.]
\item[\textbf{P2.}] \emph{The scheduling fingerprint.} Prelabor cesareans bunch
at $\arg\min_t w(t)$ (weekday business hours) and fall on weekends and
holidays; in-labor cesareans inherit $\tilde{t}$'s uniform timing and weakly
\emph{rise} on weekends (marginal-$\theta$ women who labor before being
scheduled convert). \hfill [\emph{Tested:} day-of-week, holiday and
hour-of-birth clustering, and the prelabor/in-labor split: Figures
\ref{fig:dow}--\ref{fig:hours}, Tables \ref{tab:scheduling} and
\ref{tab:mechanism_checks}.]
\item[\textbf{P3.}] \emph{Scarcity.} $\Delta$ and $\kappa$ rise where
obstetricians per birth are scarce (fewer substitutes to cover on-call time), so
weekend dips are larger there. \hfill [\emph{Tested:} density interaction, Table
\ref{tab:heterogeneity}.]
\item[\textbf{P4.}] \emph{The early-term cost.} Guaranteeing $t < \tilde{t}$
requires booking before the week-39 hazard spike, shifting prelabor cesareans to
37--38 weeks and producing excess early-term birth. \hfill [\emph{Tested:}
gestational-age distributions and early-term gap, Figure \ref{fig:gestation},
Table \ref{tab:health}.]
\item[\textbf{P5.}] \emph{Supply, not demand.} If the timing preference were the
mother's, the weekend dip would concentrate among mothers with the strongest
demand for scheduling (proxied by schooling); it does not. \hfill
[\emph{Tested:} education interaction, Table \ref{tab:heterogeneity}.]
\end{itemize}

\subsection{Policy corollary}

Interventions on $f$, such as equalizing or even tripling the vaginal fee, leave
$\theta^{*}$ nearly unchanged because the operative wedge is $\Pi$, not the fee
gap. The model thereby rationalizes our fee nulls jointly with the timing
evidence, and directs attention to organizational levers that reduce $\Pi$
itself: on-call/laborist shift coverage (which lowers $\kappa$ and detaches
$w(\tilde t)$ from the attending physician's calendar), midwife-led continuity
of care, and hospital scheduling norms.
)
%   Appendix B: data and variable construction
% =============================================================================

\clearpage
\phantomsection
\addcontentsline{toc}{section}{Appendix}
\begin{center}{\Large\bfseries Appendix}\end{center}
\vspace{0.5em}

% =============================================================================
% model.tex — A simple model of delivery-mode choice and physician time.
% % =============================================================================
% model.tex — A simple model of delivery-mode choice and physician time.
% % =============================================================================
% model.tex — A simple model of delivery-mode choice and physician time.
% \input{model} from paper.tex. Delivers the five predictions tested in the
% empirical sections (P1-P5) and the policy corollary.
% =============================================================================

\section{A simple model of delivery-mode choice and physician time}
\label{sec:model}

\subsection{Setup}

A physician attends a pregnancy whose labor, absent intervention, begins at a
random time $\tilde{t}$ distributed (approximately) uniformly over the days and
hours of the week. Clinically, the onset of spontaneous labor does not respect
the calendar, which is the identifying assumption behind our timing tests. The
physician chooses between attending a vaginal delivery ($v$) and performing a
\emph{prelabor scheduled} cesarean ($c$) at a time $t$ of her choosing; if labor
starts before a cesarean is performed, the delivery may still convert to an
\emph{in-labor} cesarean for clinical reasons.

\paragraph{Fees.} The physician receives $f_c$ for a cesarean and $f_v$ for a
vaginal delivery, where $f_v$ includes the separately-billed hourly labor
assistance, the \emph{economic} vaginal fee. In our data $f_c - f_v \le 0$ in
the largest states.

\paragraph{Time.} A cesarean takes a fixed block $\tau_c$ ($\approx$1 hour) at
the chosen time $t$. A vaginal delivery requires attending labor for a random
duration $\tau_v$ (billed mean $\approx$2.7 hours, with $E[\tau_v] > \tau_c$) at
the random time $\tilde{t}$. The opportunity cost of physician time is
\[
w(t) \;=\;
\begin{cases}
w_0 & t \in \text{weekday business hours}, \\
w_0 + \Delta & t \in \text{nights, weekends, holidays},
\end{cases}
\qquad \Delta > 0 .
\]
Committing to an event of unknown timing additionally carries an on-call cost
$\kappa \ge 0$: the physician must hold open calendar slots she could otherwise
fill (an option cost of unschedulable time).

\paragraph{Agency.} A vaginal delivery yields the patient a net clinical benefit
$\theta \sim F$ (with $\theta < 0$ when a cesarean is medically indicated). The
physician is an imperfect agent and internalizes $\alpha\theta$, $\alpha \in
(0,1]$.

\subsection{The delivery-mode decision}

The physician schedules a prelabor cesarean if and only if
\begin{equation}
\underbrace{f_c - f_v}_{\text{fee gap } \le 0}
\;+\;
\underbrace{\Pi}_{\text{convenience premium}}
\;>\; \alpha\,\theta,
\qquad
\Pi \;\equiv\; E\!\left[w(\tilde{t})\right] E[\tau_v] + \kappa - w_0\,\tau_c \;>\; 0,
\label{eq:choice}
\end{equation}
so the cutoff is $\theta^{*} = \left(f_c - f_v + \Pi\right)/\alpha$ and the
cesarean share is $F(\theta^{*})$. The convenience premium $\Pi$ is strictly
positive even when the fee gap is negative: labor falls disproportionately on
high-$w$ times ($E[w(\tilde{t})] > w_0$), lasts longer ($E[\tau_v] > \tau_c$),
and blocks the calendar ($\kappa$). Conditional on choosing $c$, the physician
schedules it at $\arg\min_t w(t)$ (weekday business hours) and, to
guarantee the procedure precedes $\tilde{t}$ (whose hazard rises steeply after
gestational week 39), books it at weeks 37--39, earlier than spontaneous labor
would deliver.

\subsection{Predictions and the corresponding tests}

\begin{itemize}
\item[\textbf{P1.}] \emph{Fees are not the operative margin.} The cesarean share
responds to the fee gap only through $f(\theta^{*})/\alpha$; when $\Pi$
dominates the wedge in \eqref{eq:choice}, the discretionary mass is already
inframarginal and fee (semi-)elasticities are $\approx 0$.
\hfill [\emph{Tested:} fee-gap regressions and fee-swing falsification, Tables
\ref{tab:not_a_price} and \ref{tab:fee_shock}; even a court order to triple
vaginal pay changed nothing.]
\item[\textbf{P2.}] \emph{The scheduling fingerprint.} Prelabor cesareans bunch
at $\arg\min_t w(t)$ (weekday business hours) and fall on weekends and
holidays; in-labor cesareans inherit $\tilde{t}$'s uniform timing and weakly
\emph{rise} on weekends (marginal-$\theta$ women who labor before being
scheduled convert). \hfill [\emph{Tested:} day-of-week, holiday and
hour-of-birth clustering, and the prelabor/in-labor split: Figures
\ref{fig:dow}--\ref{fig:hours}, Tables \ref{tab:scheduling} and
\ref{tab:mechanism_checks}.]
\item[\textbf{P3.}] \emph{Scarcity.} $\Delta$ and $\kappa$ rise where
obstetricians per birth are scarce (fewer substitutes to cover on-call time), so
weekend dips are larger there. \hfill [\emph{Tested:} density interaction, Table
\ref{tab:heterogeneity}.]
\item[\textbf{P4.}] \emph{The early-term cost.} Guaranteeing $t < \tilde{t}$
requires booking before the week-39 hazard spike, shifting prelabor cesareans to
37--38 weeks and producing excess early-term birth. \hfill [\emph{Tested:}
gestational-age distributions and early-term gap, Figure \ref{fig:gestation},
Table \ref{tab:health}.]
\item[\textbf{P5.}] \emph{Supply, not demand.} If the timing preference were the
mother's, the weekend dip would concentrate among mothers with the strongest
demand for scheduling (proxied by schooling); it does not. \hfill
[\emph{Tested:} education interaction, Table \ref{tab:heterogeneity}.]
\end{itemize}

\subsection{Policy corollary}

Interventions on $f$, such as equalizing or even tripling the vaginal fee, leave
$\theta^{*}$ nearly unchanged because the operative wedge is $\Pi$, not the fee
gap. The model thereby rationalizes our fee nulls jointly with the timing
evidence, and directs attention to organizational levers that reduce $\Pi$
itself: on-call/laborist shift coverage (which lowers $\kappa$ and detaches
$w(\tilde t)$ from the attending physician's calendar), midwife-led continuity
of care, and hospital scheduling norms.
 from paper.tex. Delivers the five predictions tested in the
% empirical sections (P1-P5) and the policy corollary.
% =============================================================================

\section{A simple model of delivery-mode choice and physician time}
\label{sec:model}

\subsection{Setup}

A physician attends a pregnancy whose labor, absent intervention, begins at a
random time $\tilde{t}$ distributed (approximately) uniformly over the days and
hours of the week. Clinically, the onset of spontaneous labor does not respect
the calendar, which is the identifying assumption behind our timing tests. The
physician chooses between attending a vaginal delivery ($v$) and performing a
\emph{prelabor scheduled} cesarean ($c$) at a time $t$ of her choosing; if labor
starts before a cesarean is performed, the delivery may still convert to an
\emph{in-labor} cesarean for clinical reasons.

\paragraph{Fees.} The physician receives $f_c$ for a cesarean and $f_v$ for a
vaginal delivery, where $f_v$ includes the separately-billed hourly labor
assistance, the \emph{economic} vaginal fee. In our data $f_c - f_v \le 0$ in
the largest states.

\paragraph{Time.} A cesarean takes a fixed block $\tau_c$ ($\approx$1 hour) at
the chosen time $t$. A vaginal delivery requires attending labor for a random
duration $\tau_v$ (billed mean $\approx$2.7 hours, with $E[\tau_v] > \tau_c$) at
the random time $\tilde{t}$. The opportunity cost of physician time is
\[
w(t) \;=\;
\begin{cases}
w_0 & t \in \text{weekday business hours}, \\
w_0 + \Delta & t \in \text{nights, weekends, holidays},
\end{cases}
\qquad \Delta > 0 .
\]
Committing to an event of unknown timing additionally carries an on-call cost
$\kappa \ge 0$: the physician must hold open calendar slots she could otherwise
fill (an option cost of unschedulable time).

\paragraph{Agency.} A vaginal delivery yields the patient a net clinical benefit
$\theta \sim F$ (with $\theta < 0$ when a cesarean is medically indicated). The
physician is an imperfect agent and internalizes $\alpha\theta$, $\alpha \in
(0,1]$.

\subsection{The delivery-mode decision}

The physician schedules a prelabor cesarean if and only if
\begin{equation}
\underbrace{f_c - f_v}_{\text{fee gap } \le 0}
\;+\;
\underbrace{\Pi}_{\text{convenience premium}}
\;>\; \alpha\,\theta,
\qquad
\Pi \;\equiv\; E\!\left[w(\tilde{t})\right] E[\tau_v] + \kappa - w_0\,\tau_c \;>\; 0,
\label{eq:choice}
\end{equation}
so the cutoff is $\theta^{*} = \left(f_c - f_v + \Pi\right)/\alpha$ and the
cesarean share is $F(\theta^{*})$. The convenience premium $\Pi$ is strictly
positive even when the fee gap is negative: labor falls disproportionately on
high-$w$ times ($E[w(\tilde{t})] > w_0$), lasts longer ($E[\tau_v] > \tau_c$),
and blocks the calendar ($\kappa$). Conditional on choosing $c$, the physician
schedules it at $\arg\min_t w(t)$ (weekday business hours) and, to
guarantee the procedure precedes $\tilde{t}$ (whose hazard rises steeply after
gestational week 39), books it at weeks 37--39, earlier than spontaneous labor
would deliver.

\subsection{Predictions and the corresponding tests}

\begin{itemize}
\item[\textbf{P1.}] \emph{Fees are not the operative margin.} The cesarean share
responds to the fee gap only through $f(\theta^{*})/\alpha$; when $\Pi$
dominates the wedge in \eqref{eq:choice}, the discretionary mass is already
inframarginal and fee (semi-)elasticities are $\approx 0$.
\hfill [\emph{Tested:} fee-gap regressions and fee-swing falsification, Tables
\ref{tab:not_a_price} and \ref{tab:fee_shock}; even a court order to triple
vaginal pay changed nothing.]
\item[\textbf{P2.}] \emph{The scheduling fingerprint.} Prelabor cesareans bunch
at $\arg\min_t w(t)$ (weekday business hours) and fall on weekends and
holidays; in-labor cesareans inherit $\tilde{t}$'s uniform timing and weakly
\emph{rise} on weekends (marginal-$\theta$ women who labor before being
scheduled convert). \hfill [\emph{Tested:} day-of-week, holiday and
hour-of-birth clustering, and the prelabor/in-labor split: Figures
\ref{fig:dow}--\ref{fig:hours}, Tables \ref{tab:scheduling} and
\ref{tab:mechanism_checks}.]
\item[\textbf{P3.}] \emph{Scarcity.} $\Delta$ and $\kappa$ rise where
obstetricians per birth are scarce (fewer substitutes to cover on-call time), so
weekend dips are larger there. \hfill [\emph{Tested:} density interaction, Table
\ref{tab:heterogeneity}.]
\item[\textbf{P4.}] \emph{The early-term cost.} Guaranteeing $t < \tilde{t}$
requires booking before the week-39 hazard spike, shifting prelabor cesareans to
37--38 weeks and producing excess early-term birth. \hfill [\emph{Tested:}
gestational-age distributions and early-term gap, Figure \ref{fig:gestation},
Table \ref{tab:health}.]
\item[\textbf{P5.}] \emph{Supply, not demand.} If the timing preference were the
mother's, the weekend dip would concentrate among mothers with the strongest
demand for scheduling (proxied by schooling); it does not. \hfill
[\emph{Tested:} education interaction, Table \ref{tab:heterogeneity}.]
\end{itemize}

\subsection{Policy corollary}

Interventions on $f$, such as equalizing or even tripling the vaginal fee, leave
$\theta^{*}$ nearly unchanged because the operative wedge is $\Pi$, not the fee
gap. The model thereby rationalizes our fee nulls jointly with the timing
evidence, and directs attention to organizational levers that reduce $\Pi$
itself: on-call/laborist shift coverage (which lowers $\kappa$ and detaches
$w(\tilde t)$ from the attending physician's calendar), midwife-led continuity
of care, and hospital scheduling norms.
 from paper.tex. Delivers the five predictions tested in the
% empirical sections (P1-P5) and the policy corollary.
% =============================================================================

\section{A simple model of delivery-mode choice and physician time}
\label{sec:model}

\subsection{Setup}

A physician attends a pregnancy whose labor, absent intervention, begins at a
random time $\tilde{t}$ distributed (approximately) uniformly over the days and
hours of the week. Clinically, the onset of spontaneous labor does not respect
the calendar, which is the identifying assumption behind our timing tests. The
physician chooses between attending a vaginal delivery ($v$) and performing a
\emph{prelabor scheduled} cesarean ($c$) at a time $t$ of her choosing; if labor
starts before a cesarean is performed, the delivery may still convert to an
\emph{in-labor} cesarean for clinical reasons.

\paragraph{Fees.} The physician receives $f_c$ for a cesarean and $f_v$ for a
vaginal delivery, where $f_v$ includes the separately-billed hourly labor
assistance, the \emph{economic} vaginal fee. In our data $f_c - f_v \le 0$ in
the largest states.

\paragraph{Time.} A cesarean takes a fixed block $\tau_c$ ($\approx$1 hour) at
the chosen time $t$. A vaginal delivery requires attending labor for a random
duration $\tau_v$ (billed mean $\approx$2.7 hours, with $E[\tau_v] > \tau_c$) at
the random time $\tilde{t}$. The opportunity cost of physician time is
\[
w(t) \;=\;
\begin{cases}
w_0 & t \in \text{weekday business hours}, \\
w_0 + \Delta & t \in \text{nights, weekends, holidays},
\end{cases}
\qquad \Delta > 0 .
\]
Committing to an event of unknown timing additionally carries an on-call cost
$\kappa \ge 0$: the physician must hold open calendar slots she could otherwise
fill (an option cost of unschedulable time).

\paragraph{Agency.} A vaginal delivery yields the patient a net clinical benefit
$\theta \sim F$ (with $\theta < 0$ when a cesarean is medically indicated). The
physician is an imperfect agent and internalizes $\alpha\theta$, $\alpha \in
(0,1]$.

\subsection{The delivery-mode decision}

The physician schedules a prelabor cesarean if and only if
\begin{equation}
\underbrace{f_c - f_v}_{\text{fee gap } \le 0}
\;+\;
\underbrace{\Pi}_{\text{convenience premium}}
\;>\; \alpha\,\theta,
\qquad
\Pi \;\equiv\; E\!\left[w(\tilde{t})\right] E[\tau_v] + \kappa - w_0\,\tau_c \;>\; 0,
\label{eq:choice}
\end{equation}
so the cutoff is $\theta^{*} = \left(f_c - f_v + \Pi\right)/\alpha$ and the
cesarean share is $F(\theta^{*})$. The convenience premium $\Pi$ is strictly
positive even when the fee gap is negative: labor falls disproportionately on
high-$w$ times ($E[w(\tilde{t})] > w_0$), lasts longer ($E[\tau_v] > \tau_c$),
and blocks the calendar ($\kappa$). Conditional on choosing $c$, the physician
schedules it at $\arg\min_t w(t)$ (weekday business hours) and, to
guarantee the procedure precedes $\tilde{t}$ (whose hazard rises steeply after
gestational week 39), books it at weeks 37--39, earlier than spontaneous labor
would deliver.

\subsection{Predictions and the corresponding tests}

\begin{itemize}
\item[\textbf{P1.}] \emph{Fees are not the operative margin.} The cesarean share
responds to the fee gap only through $f(\theta^{*})/\alpha$; when $\Pi$
dominates the wedge in \eqref{eq:choice}, the discretionary mass is already
inframarginal and fee (semi-)elasticities are $\approx 0$.
\hfill [\emph{Tested:} fee-gap regressions and fee-swing falsification, Tables
\ref{tab:not_a_price} and \ref{tab:fee_shock}; even a court order to triple
vaginal pay changed nothing.]
\item[\textbf{P2.}] \emph{The scheduling fingerprint.} Prelabor cesareans bunch
at $\arg\min_t w(t)$ (weekday business hours) and fall on weekends and
holidays; in-labor cesareans inherit $\tilde{t}$'s uniform timing and weakly
\emph{rise} on weekends (marginal-$\theta$ women who labor before being
scheduled convert). \hfill [\emph{Tested:} day-of-week, holiday and
hour-of-birth clustering, and the prelabor/in-labor split: Figures
\ref{fig:dow}--\ref{fig:hours}, Tables \ref{tab:scheduling} and
\ref{tab:mechanism_checks}.]
\item[\textbf{P3.}] \emph{Scarcity.} $\Delta$ and $\kappa$ rise where
obstetricians per birth are scarce (fewer substitutes to cover on-call time), so
weekend dips are larger there. \hfill [\emph{Tested:} density interaction, Table
\ref{tab:heterogeneity}.]
\item[\textbf{P4.}] \emph{The early-term cost.} Guaranteeing $t < \tilde{t}$
requires booking before the week-39 hazard spike, shifting prelabor cesareans to
37--38 weeks and producing excess early-term birth. \hfill [\emph{Tested:}
gestational-age distributions and early-term gap, Figure \ref{fig:gestation},
Table \ref{tab:health}.]
\item[\textbf{P5.}] \emph{Supply, not demand.} If the timing preference were the
mother's, the weekend dip would concentrate among mothers with the strongest
demand for scheduling (proxied by schooling); it does not. \hfill
[\emph{Tested:} education interaction, Table \ref{tab:heterogeneity}.]
\end{itemize}

\subsection{Policy corollary}

Interventions on $f$, such as equalizing or even tripling the vaginal fee, leave
$\theta^{*}$ nearly unchanged because the operative wedge is $\Pi$, not the fee
gap. The model thereby rationalizes our fee nulls jointly with the timing
evidence, and directs attention to organizational levers that reduce $\Pi$
itself: on-call/laborist shift coverage (which lowers $\kappa$ and detaches
$w(\tilde t)$ from the attending physician's calendar), midwife-led continuity
of care, and hospital scheduling norms.
   % Appendix A: the model

% =============================================================================
\section{Data and variable construction}\label{app:data}
% =============================================================================
This appendix documents the construction of the analysis files; the complete
codebook is distributed in machine-readable form with the replication package
described in the data-availability statement. Deliveries in the TISS claims are identified by
their procedure codes in the Unified Health Supplementary Terminology
(Terminologia Unificada da Sa\'ude Suplementar, TUSS), the standardized procedure
nomenclature that private plans use to bill (cesarean 31309054 and 31309208;
vaginal 31309127); the economic vaginal fee adds the separately billed hourly
labor-assistance code (31309038). In SINASC, the sector of a birth is the
legal-entity type of its establishment in the CNES registry: for-profit (codes
2xxx), nonprofit (3xxx), and public (1xxx); the headline contrast
compares for-profit with public establishments. All files are
built at the municipality level on the six-digit municipality code of the
Brazilian Institute of Geography and Statistics (Instituto Brasileiro de
Geografia e Estat\'istica, IBGE).

\paragraph{Variable definitions.} The main constructed variables are defined as
follows. Indicators equal one when the stated condition holds and zero otherwise.
\begin{itemize}
\item \textbf{Cesarean}: the delivery is a cesarean section. In TISS this is
the presence of a cesarean procedure code; in SINASC it is the registry's
delivery-type field.
\item \textbf{Sector (for-profit / nonprofit / public)}: the sector of the birth
establishment, from the legal-entity type of its CNES record; \textbf{for-profit}
denotes establishments with legal-entity type 2xxx, the headline group.
\item \textbf{Economic vaginal fee}: the physician fee for a vaginal delivery
plus the separately billed hourly labor-assistance fee (TUSS 31309038), which
compensates the time spent attending labor. This is the correct comparison for
the cesarean fee because it credits the vaginal delivery with the time it pays for.
\item \textbf{Log fee gap}: the natural logarithm of the ratio of the mean
cesarean fee to the mean economic vaginal fee in a municipality-year; positive
values mean cesareans pay more.
\item \textbf{Cesarean share}: the share of births delivered by cesarean in a
municipality-date-sector cell.
\item \textbf{Weekend, holiday, eve}: indicators for Saturdays and Sundays, the
eight fixed and four Easter-based movable national holidays, and the eve of a rest
day (a Friday or the day before a holiday), respectively.
\item \textbf{Robson group}: the World Health Organization's ten-group
classification of deliveries; groups 1 and 2 are the low-risk first births used in
the mechanism tests.
\item \textbf{Prelabor cesarean}: a cesarean recorded as performed before labor
began, as opposed to an in-labor cesarean.
\item \textbf{Early-term birth}: a birth at 37 or 38 completed weeks of
gestation.
\item \textbf{Treated (Parto Adequado)}: the municipality contains a private
hospital that joined the program's 2017--2021 dissemination phase.
\item \textbf{Obstetricians per 1,000 births}: the count of obstetricians in
the CNES professionals file per 1,000 births in the municipality-year.
\end{itemize}
